% 原书第 21—22 页；第 2.3 节导言及第 2.3.1—2.3.2 小节。
\section{定义因果目标参数}
\label{sec:ch02-causal-target}

有了描述数据生成机制（data-generating mechanism）的 SCM，现在我们进一步定义因果目标参数。确切地说，我们将其定义为 SCM 中完整数据（full data）$(U,X)$ 的分布的一个参数。形式上，将完整数据 $(U,X)$ 的 SCM 记为 $\mathcal M^F$，它是 SCM 所描述的一组可能分布 $P_{U,X}$ 的集合。作为完整数据的模型，$\mathcal M^F$ 收集了底层数据 $(U,X)$ 的所有可能分布。$\Psi^F$ 是作用于 $P_{U,X}$ 的映射，所得 $\Psi^F(P_{U,X})$ 就是 $P_{U,X}$ 的目标参数。对于在所设定 SCM 编码的假设下，$(U,X)$ 的每一个可能分布 $P_{U,X}$，都需要定义这一映射。因此，我们写为 $\Psi^F\colon\mathcal M^F\to\mathbb R^d$，其中，$\mathbb R^d$ 表示参数是由 $d$ 个实数组成的向量。SCM $\mathcal M^F$ 包含的分布，由确定性函数 $f=(f_{X_j}:j)$ 和 $U$ 的分布 $P_U$ 索引；而 $f$ 和这一联合分布 $P_U$，都可以从完整数据 $(U,X)$ 的分布中识别。因此，目标参数也可以表示为 $f$ 与 $U$ 的联合分布的函数。

对于 IVF 研究，观测数据为 $O=(L_0,A_0,L_1,A_1,Y)$，SCM 由式 (2.1) 给出，且不对分布 $P_U$ 作任何假设。我们可以定义随机变量 $Y_{a_{0:1}}=f_Y(L_0,a_0,L_1,a_1,U_Y)$，对应于 SCM 中的干预 $A_{0:1}=a_{0:1}$。$Y_{a_{0:1}}$ 的边缘概率分布为
\[
 P_{U,X}(Y_{a_{0:1}}=y)
 =P_{U,X}\bigl(f_Y(L_0,a_0,L_1,a_1,U_Y)=y\bigr).
\]
回顾第 1 章给出的感兴趣统计参数：
% 本式与第 1 章原式一致：求和号不补下标，L_1=1 不改为 L_1=l_1。
\[
 \begin{aligned}
 \Psi(P_0)=E_{P_0}\Bigl(&\sum P_0(Y=1\mid A_{0:1}=1,L_{0:1}=l_{0:1})\\
 &{}\times P_0(L_1=1\mid A_0=1,L_0=l_0)\Bigr).
 \end{aligned}
\]
后文将说明，在随机化假设和正值性（positivity）假设下，这一统计参数等于 $(U,X)$ 分布的因果目标参数：
\[
 \Psi^F(P_{U,X})=P_{U,X}(Y_{(1,1)}=1).
\]

\subsection{干预}
\label{subsec:ch02-interventions}

我们可以将干预节点 $A_{0:1}$ 设为某些取值 $a_{0:1}\in\mathcal A$，从而干预 SCM 定义的系统。其中，$\mathcal A$ 是暴露（exposure），即 IVF 治疗，所有可能取值的集合。通过干预，我们能够描述系统在干预变量的指定水平下本会生成的数据。这是一个至关重要的概念，因为我们将因果目标参数定义为：在对 $f$ 中结构方程实施干预后，数据 $(U,X)$ 的分布的一个参数。干预定义了一个作为 $(U,X)$ 函数的随机变量，因而，目标参数为 $\Psi^F(P_{U,X})$。

在 IVF 研究中，我们可以对 $A_{0:1}$ 进行干预，以考察特定 IVF 治疗暴露下本会发生什么。需要强调的是，对 SCM 实施干预，即改变干预变量对应的函数 $f_{X_j}$，并不会改变 $f$ 中的其他函数。对于式 (2.1) 的 SCM，我们可以干预 $f_{A_0}$ 和 $f_{A_1}$，将 $a_0$ 和 $a_1$ 都设为 1：
\[
 \begin{aligned}
 L_0&=f_{L_0}(U_{L_0}),\\
 A_0&=1,\\
 L_1&=f_{L_1}(L_0,1,U_{L_1}),\\
 A_1&=1,\\
 Y_{(1,1)}&=f_Y(L_0,1,L_1,1,U_Y).
 \end{aligned}
\]
这一干预定义了一个作为 $(U,X)$ 函数的随机变量，即 $a_0=1$、$a_1=1$ 时的 $Y_{a_{0:1}}=Y_{a_{0:1}}(U)$。我们的目标参数是干预后分布（postintervention distribution）的一个参数；这一分布，就是干预下 $(X,U)$ 的概率分布。因此，完整数据的 SCM 使我们能够对每个 $a_{0:1}$，定义随机变量 $Y_{a_{0:1}}=f_Y(L_0,a_0,L_1,a_1,U_Y)$；这里，$Y_{a_{0:1}}$ 表示在这一系统下，某个特定研究对象接受暴露 $a_{0:1}$ 时本会观测到的结局。

\subsection{反事实}
\label{subsec:ch02-counterfactuals}

在所有可能的暴露轨迹下都观测每名研究对象的“理想实验”是不可能实现的。每名研究对象只能贡献一个 $Y$，即在其实际经历的暴露下观测到的结局。上面，我们通过干预 $A_{0:1}$，将 $a_0=1$ 且 $a_1=1$，从而生成每名研究对象在接受两轮 IVF 治疗条件下的结局。回顾前面的定义，$Y_{a_{0:1}}$ 表示这一系统下某个特定研究对象接受暴露 $a_{0:1}$ 时本会观测到的结局；我们有 $(Y_{a_{0:1}}:a_{0:1})$，其中 $a_{0:1}\in\mathcal A$。对于研究中的每个实现值 $u$，也就是从目标总体随机抽取的一名个体，我们通过干预式 (2.1)，生成反事实结局（counterfactual outcome）$Y_{(1,1)}(u)$。反事实结局由 SCM 所蕴含，是 SCM 的推论。也就是说，$Y_{(1,1)}(u)=f_Y(L_0,1,L_1,1,u_Y)$。随机反事实 $Y_{(1,1)}(U)$ 的随机性，来自 $U$ 的概率分布。
